%% kapa: exam=1 round=33 year=2022 period=1 form=A subject=economics subject_ko=경제학원론 seq=18 number=58
%% source: Q-Net 감정평가사 기출문제 (cst00309) artlSeq=5209997 file=raw/1008/5209997/2022년도 제33회 감정평가사 1차시험 1교시(민법,경제학,부동산학).pdf
%% conversion: kordoc 4.24.51 -> tools/gen1.py (auto); review=proofread-sonnet; flags=none
\begin{question}{58}
A국에서는 교역 이전 X재의 국내가격이 국제가격보다 더 높다. 교역 이후 국제가격으로 A국이 X재의 초과수요분을 수입한다면, 이로 인해 A국에 나타나는 효과로 옳은 것은? (단, 공급곡선은 우상향, 수요곡선은 우하향한다.)
\begin{choices}
  \item 교역 전과 비교하여 교역 후 생산자 잉여가 감소한다.
  \item 교역 전과 비교하여 교역 후 소비자 잉여가 감소한다.
  \item 생산자 잉여는 교역 여부와 무관하게 일정하다.
  \item 교역 전과 비교하여 교역 후 총잉여가 감소한다.
  \item 총잉여는 교역 여부와 무관하게 일정하다.
\end{choices}
\end{question}
\answer{1}
